\section{The Brown Scapular: Received Sign and Disciplined Tradition}

\subsection{What the object is}

The Brown Scapular is a reduced form of the Carmelite religious scapular: cloth falling before and behind from the shoulders as part of the habit. For lay use it is ordinarily two small pieces of brown or dark cloth joined by cords. Its sacramental meaning comes from the Church's blessing, Carmelite affiliation, the wearer's faith and disposition, and the Paschal mystery to which all sacramentals point---not from fabric acting by an inherent occult force.

The Catechism, nos.~1667--1670, distinguishes sacramentals from sacraments. Sacraments confer grace by Christ's institution; sacramentals are sacred signs instituted by the Church that prepare the faithful to receive grace and dispose them to cooperate with it. The Scapular does not baptize, absolve, communicate the Eucharist, or guarantee final perseverance. It recalls and supports the life already begun in Baptism.

The Holy See's \textit{Directory}, no.~205, gives the controlling synthesis. The Scapular is the reduced Carmelite habit, an external sign of filial relation to Mary, trust in her maternal intercession, the primacy of spiritual life, and the need for prayer. The approved rite describes it as a remembrance that in Baptism the Christian has been clothed in Christ. Its imposition should follow preparation and should not be improvised.

\subsection{The 1251 tradition and its documentary limit}

Carmelite tradition associates the Scapular with an appearance of Mary to an English Carmelite, Simon Stock, conventionally dated 16 July 1251 and often located at Aylesford. The familiar promise says, in varying forms, that one who dies wearing the habit will not suffer eternal fire. John Paul II in 2001 called the bestowal a venerable tradition of the Order; he did not issue a modern judgment declaring the event historically or supernaturally proved.

The surviving evidence is late. Secure thirteenth-century documentation about Simon is sparse; several medieval Simons may have been combined; the developed vision story becomes visible only near the end of the fourteenth century; and a Carmelite participant in a 1375 Cambridge dispute about Marian bestowal of religious habits appears not to know it. Richard Copsey's documentary study and current Carmelite historical catechesis therefore judge the provenance disputed.

The proper conclusion is neither contempt nor credulity. The story may be narrated reverently as received Carmelite tradition. It may not be promoted as a defined fact, used to accuse historical inquiry of impiety, or treated as the foundation without which the approved sacramental collapses. The local controlling record links this booklet to the repository's Marian-apparition status dossier, which locates no positive event decree for the 1251 claim.

\subsection{How to read the promise}

No private revelation can repeal the Gospel's doctrine of judgment, conversion, grace, and perseverance. ``Whoever perseveres to the end will be saved'' (Mt 24:13); the crown of life is promised to fidelity unto death (Rev 2:10). The most defensible meaning of the Scapular promise is therefore ecclesial and dispositional: Mary intercedes for those who truly live and die in Christ; the habit-sign calls the wearer to the Carmelite form of such fidelity.

Simply retaining cloth while freely rejecting God cannot turn mortal sin into communion. Conversely, losing the cloth by accident does not overpower divine mercy. Final perseverance is grace; the sacraments are the ordinary means Christ entrusted to the Church; repentance remains necessary; and God's judgment is neither manipulated nor deceived.

\subsection{The Sabbatine claim}

A later tradition, attached to a bull attributed to John XXII in 1322, promised release from purgatory on the first Saturday after death under specified conditions. Historical criticism does not accept that bull as authentic. The Holy Office's 1613 intervention did not make a forged document genuine or authorize a mechanical timetable. It allowed a carefully bounded confidence that Mary, by continual intercession and special protection, assists deceased Carmelite brothers, sisters, and confraternity members who died in charity and lived the stated Christian commitments, with Saturday named as her devotional day.

Current Carmelite catechesis does not promulgate a literal ``first Saturday release'' as the doctrine of the Scapular. The sound remainder is ordinary Catholic teaching: the faithful may pray for the dead; the saints intercede; Mary exercises maternal care; purification after death is governed by God's holy mercy; and no calendar formula binds the sovereign Judge. This booklet therefore omits the Sabbatine conditions from its prayer rule and makes no purgatorial timetable part of the novena.

\subsection{Blessing, enrollment, and daily life}

The Congregation for Divine Worship and the Discipline of the Sacraments approved the present Rite of Blessing and Enrollment on 5 January 1996. Current Carmelite catechesis states that a priest or deacon may bless and impose the Scapular; reception into a confraternity follows the rite and competent Carmelite authority. A lay person should not invent a self-enrollment or assume that buying and wearing cloth creates a juridic status. Local Carmelite or parish guidance should be sought where enrollment is desired.

Enrollment is not required to pray this novena. Nor does praying the novena enroll anyone. A person may honor Mary and study Carmelite spirituality without wearing the Scapular. One who does wear it should receive its language as a vocation:

\begin{itemize}
\item frequent participation in Mass and worthy reception of Holy Communion;
\item regular meditation on Sacred Scripture and, according to capacity, the Liturgy of the Hours;
\item prayer and Marian imitation ordered to Christ;
\item chastity according to one's state, charity, obedience to God's will, and works of mercy;
\item real affiliation to the Church and some living contact with Carmelite spirituality, rather than a good-luck object.
\end{itemize}

The cloth may be simple and undecorated; pictures are not what constitute it. A medal has been permitted as a substitute in specified circumstances, but cloth better preserves the visible relation to the habit. These practical questions matter less than the governing truth: the exterior sign must become an interior and social ``habit'' of Christlike life.

\subsection{Protection without magic}

Mary's protection is real intercession within Providence, not a shield against every temporal cross. It may include rescue, consolation, correction, strength to repent, sacramental reconciliation, courage in danger, patience in illness, and assistance at death; God determines its form. A claim that the Scapular guarantees immunity from injury, fire, war, temptation, mental illness, or natural consequences exceeds the Church's teaching.

The safest visible test is transformation. Does wearing the sign make the wearer easier to correct, quicker to forgive, more faithful in prayer, more truthful, purer in conduct, more attentive to the poor, and more ready for the sacraments? If not, the remedy is not a more elaborate object but conversion.
